% BHU PYQ -- Professional Exam Template
\documentclass[11pt,a4paper]{article}

\usepackage[a4paper,top=2.5cm,bottom=2.5cm,left=2.8cm,right=2.8cm,headheight=14pt]{geometry}
\usepackage{amsmath,amssymb}
\usepackage[T1]{fontenc}
\usepackage[utf8]{inputenc}
\usepackage{lmodern}
\usepackage{microtype}
\usepackage{fancyhdr}
\usepackage{enumitem}
\usepackage{listings}
\usepackage{hyperref}
\usepackage{lastpage}
\usepackage{parskip}

\hypersetup{colorlinks=true,urlcolor=black,linkcolor=black,
  pdftitle={CS-101: Problem Solving through C Programming -- BHU PYQ 2023-24}}

\pagestyle{fancy}
\fancyhf{}
\renewcommand{\headrulewidth}{0.4pt}
\renewcommand{\footrulewidth}{0.4pt}
\fancyhead[L]{\small   \textit{Banaras Hindu University}}
\fancyhead[R]{\small   \textit{CS-101: Problem Solving through C Programming}}
\fancyfoot[C]{\small Page  \thepage\ of \pageref{LastPage}}

\lstset{
  basicstyle=\small tfamily,
  frame=single,
  breaklines=true,
  columns=flexible,
  keepspaces=true
}


\newlist{parts}{enumerate}{2}
\setlist[parts,1]{label=(\alph*),leftmargin=*,itemsep=4pt,topsep=3pt}
\setlist[parts,2]{label=(\roman*),leftmargin=*,itemsep=2pt,topsep=2pt}

\newcommand{\pts}[1]{\hfill   \textit{[#1]}}


\begin{document}

\begin{center}
  {\large\bfseries BANARAS HINDU UNIVERSITY}\\[2pt]
  {\normalsize B.Sc.\ (Hons.) Semester I Examination, 2023-24}\\[6pt]
  \rule{0.8\linewidth}{0.5pt}\\[6pt]
  {\large\bfseries Paper: CS-101 --- Problem Solving through C Programming}\\[4pt]
  {\normalsize Time: Three Hours \hfill Maximum Marks: 70}
\end{center}

\rule{\linewidth}{0.4pt}

\smallskip



\noindent   \textit{Instructions: Answer any   \textbf{five} questions including Question~1 which is compulsory. Write your Roll No.\ at the top immediately on receipt of this question paper.}

\bigskip

%% ─── Q1 (Compulsory) ───────────────────────────────────────────────────────



\noindent   \textbf{Question 1.}   \textit{Attempt any   \textbf{nine}.}\pts{2$ \times$9 = 18}

\begin{parts}
  \item Differentiate between   \texttt{float} and   \texttt{double} data types in C.
  \item What is the difference between the   \texttt{=} operator and the   \texttt{==} operator in C?
  \item What is the role of the   \texttt{break} keyword in C programming?
  \item What is the difference between a string and an array of characters?
  \item What is the use of the   \texttt{\#ifndef} preprocessor directive in C programming?
  \item What is the difference between   \texttt{scanf()} and   \texttt{gets()} functions?
  \item Write the syntax and usage of the   \texttt{->} operator in C.
  \item Define the concept of \emph{pointer to pointer} with a suitable example.
  \item What will be the output of the following code? (Explain with reason.)

\begin{lstlisting}[language=C]
void main() {
    int x = 10, y = 5, p, q;
    p = x - 19;
    q = (p++) && (y = 3);
    printf("p = %d  q = %d", p, q);
}
\end{lstlisting}
  \item What will be the output of the following code? (Explain with reason.)

\begin{lstlisting}[language=C]
void main() {
    int x = 3, y = 2, z;
    z = x++ + ++y;
    printf("x = %d  y = %d  z = %d", x, y, z);
}
\end{lstlisting}
\end{parts}

\bigskip

%% ─── Q2 ─────────────────────────────────────────────────────────────────────



\noindent   \textbf{Question 2.}

\begin{parts}
  \item What is an \emph{algorithm}? Draw the flowchart to print the real roots of the quadratic equation $ax^{2}+bx+c=0$ using the Sridharacharya formula.\pts{4}
  \item Explain the difference between a   \texttt{while} loop and a   \texttt{do\ldots while} loop with a suitable example.\pts{5}
  \item Explain the   \texttt{switch\ldots case} statement with a suitable example.\pts{4}
\end{parts}

\bigskip

%% ─── Q3 ─────────────────────────────────────────────────────────────────────



\noindent   \textbf{Question 3.}

\begin{parts}
  \item Discuss the \emph{logical operators} in C with a suitable example.\pts{4}
  \item Differentiate between \emph{actual} and \emph{formal} arguments when using a function. Explain \emph{call by reference} with a suitable example.\pts{5}
  \item Define \emph{recursion}. Write a program to find the factorial of a given number using recursion.\pts{4}
\end{parts}

\bigskip

%% ─── Q4 ─────────────────────────────────────────────────────────────────────



\noindent   \textbf{Question 4.}

\begin{parts}
  \item State the difference between   \texttt{register} and   \texttt{static} storage class variables with appropriate examples.\pts{4}
  \item Write your own functions in C for the following string operations:\pts{6}
    
\begin{parts}
      \item Copy one string into another.
      \item Reverse a given string.
    \end{parts}
  \item What are \emph{command line arguments}? Why are they used?\pts{3}
\end{parts}

\bigskip

%% ─── Q5 ─────────────────────────────────────────────────────────────────────



\noindent   \textbf{Question 5.}

\begin{parts}
  \item How are elements of a 2D array stored in computer memory? Give one example.\pts{2}
  \item Differentiate between the concepts of \emph{pointer to an array} and \emph{array of pointers} with suitable examples.\pts{6}
  \item How can a pointer be used to access the elements of a one-dimensional array? Explain with an example.\pts{5}
\end{parts}

\bigskip

%% ─── Q6 ─────────────────────────────────────────────────────────────────────



\noindent   \textbf{Question 6.}

\begin{parts}
  \item Explain the different possible \emph{modes} for opening a file in C. Write a program to copy the contents of one file into another.\pts{5}
  \item What do you mean by \emph{dynamic memory allocation} in C? Explain the role of the   \texttt{realloc()} function with a suitable example.\pts{4}
  \item Define \emph{structure} with a suitable example. How is memory allocation done for structure variables? Give one example.\pts{4}
\end{parts}

\bigskip

%% ─── Q7 ─────────────────────────────────────────────────────────────────────



\noindent   \textbf{Question 7.}

\begin{parts}
  \item What is the role of the \emph{Preprocessor} in C programming? Differentiate between the two forms of include directives:   \texttt{\#include<mylib.h>} and   \texttt{\#include"mylib.h"}.\pts{4}
  \item Explain the concept of \emph{nested structure} with a suitable example.\pts{5}
  \item Define \emph{union} with a suitable example. Give one real-life example where a union can be used.\pts{4}
\end{parts}

\vfill
\rule{\linewidth}{0.4pt}

\begin{center}\small
  \href{https://bhu-exam.vercel.app/}{Click here to visit our website}\\[4pt]
  \href{https://me-aryan.vercel.app/}{Click here to visit our contributor}\\[4pt]
  \href{https://quashan-me.vercel.app/}{Click here to visit our contributor}
\end{center}
\end{document}